\subsection{Выбор программных средств}

Программная система включает три компонента, для каждого из которых выбран технологический стек с учётом единства структурных решений и совместимости.

\subsubsection{Серверная часть и база данных}

В качестве СУБД выбрана PostgreSQL с расширением PostGIS~--- оно позволяет выполнять геопространственные запросы непосредственно на уровне базы данных. Функция ST\_DWithin расширения PostGIS позволяет эффективно искать обращения в заданном радиусе от указанной точки, что является основой алгоритма дедупликации.

Серверная часть реализована на C\# с использованием ASP.NET Core и Entity Framework Core. Подход Code-First позволяет описывать модель данных в виде классов C\# и автоматически генерировать миграции. Выбор обусловлен следующими преимуществами: кроссплатформенность среды выполнения .NET, встроенная поддержка HTTPS и JWT-аутентификации, интеграция с OpenAPI (Swagger) для автодокументирования API, производительность и масштабируемость. Серверная часть предоставляет единый программный интерфейс (REST API), используемый обоими клиентскими приложениями~--- веб-панелью и мобильным приложением.

\subsubsection{Клиентские приложения}

Для обоих клиентских приложений выбрана экосистема React с языком TypeScript, что обеспечивает единообразие подходов к разработке, типизации и управлению состоянием.

Диспетчерская веб-панель реализована на React с библиотекой Яндекс.Карты (JavaScript API) для интерактивных карт. Сборка выполняется инструментом Vite. Яндекс.Карты предоставляют готовый API для отображения карты, маркеров, кластеризации и геокодирования на русскоязычной картографической основе.

Мобильное приложение реализовано на React Native с инструментарием Expo. Главным доводом в пользу React Native стала кроссплатформенность (единая кодовая база для Android и iOS) и преемственность с веб-панелью: подходы к структуре компонентов, управлению состоянием и типизации остаются общими. Expo предоставляет готовые модули для доступа к камере (expo-image-picker), геолокации (expo-location), защищённого хранения токенов (expo-secure-store) и среду Expo Go для тестирования без нативной сборки. Навигация реализована через expo-router, карта~--- через react-native-maps (Google Maps на Android, Apple Maps на iOS). HTTP-взаимодействие с серверной частью в обоих клиентах выполняется через Axios с поддержкой JWT-аутентификации.
